\section{Motivic Donaldson-Thomas Theory for CHL Quotients}
\label{wn:sec:motivic-DT-CHL-N-ge-2}

\providecommand{\DTmot}{Z^{\mathrm{red},\mathrm{mot}}_{\mathrm{DT}}}
\providecommand{\FCHL}{F^{\mathrm{CHL}}}
\providecommand{\FDD}{F^{\mathrm{dd}}}
\providecommand{\Lmot}{\mathbb L^{1/2}}
\providecommand{\KSring}{\mathcal M^{\widehat\mu}_{\mathbb C}}
\providecommand{\MOmot}{\Omega^{\mathrm{mot}}}
\providecommand{\PhiOP}{\Phi_{10}^{\mathrm{OP}}}
\providecommand{\Phiun}{\Phi_{10}^{\mathrm{un}}}
\providecommand{\kBKM}{\kappa_{\mathrm{BKM}}}
\providecommand{\kcat}{\kappa_{\mathrm{cat}}}
\providecommand{\kch}{\kappa_{\mathrm{ch}}}
\providecommand{\kfib}{\kappa_{\mathrm{fiber}}}

\paragraph{The quotient.}
Let $S$ be a K3 surface, let $E$ be an elliptic curve, let $g_N$ be a
symplectic automorphism of $S$ of order $N$, and let $t_N$ be
translation by a point of exact order $N$ on $E$, where
$N\in\{1,2,3,4,6\}$. Put
\[
  G_N=\langle(g_N,t_N)\rangle,\qquad
  X_N=(S\times E)/G_N .
\]
For every $a\in\{1,\ldots,N-1\}$, the translation $t_N^a$ has no fixed
point on $E$. Hence $(g_N^a,t_N^a)$ has no fixed point on $S\times E$,
even when $g_N^a$ fixes points on $S$. The quotient map
\[
  \pi_N:S\times E\longrightarrow X_N
\]
is finite \'etale of degree $N$. The holomorphic form
$\Omega_S\wedge dz_E$ descends because $g_N^*\Omega_S=\Omega_S$ and
translations preserve $dz_E$. Thus $X_N$ is a smooth projective
threefold with trivial canonical bundle. The Du Val strata of
$S/\langle g_N\rangle$ enter as surface degeneration charts, while the
diagonal threefold quotient remains smooth.

\paragraph{Surface local data.}
The K3 quotient supplies the cyclic surface charts used in vertex and
degeneration calculations. The elliptic translation removes fixed
points from the threefold quotient and leaves the surface chart types
unchanged.

For order $4$, Nikulin-Mukai fixed point counts give four points fixed
by $g_4$ and eight points fixed by $g_4^2$. The four $g_4$-fixed points
give four $A_3$ singularities on $S/\langle g_4\rangle$. The remaining
four $g_4^2$-fixed points form two $g_4$-orbits with stabilizer
$\mathbb Z/2$, hence two $A_1$ singularities. The surface type is
\[
  4A_3+2A_1,
\]
and the local motivic factor is
\[
  \bigl(V^{\mathrm{mot}}_{A_3}\bigr)^4
  \bigl(V^{\mathrm{mot}}_{A_1}\bigr)^2 .
\]
For order $6$, the fixed point counts are
\[
  \#\mathrm{Fix}(g_6)=2,\qquad
  \#\mathrm{Fix}(g_6^2)=6,\qquad
  \#\mathrm{Fix}(g_6^3)=8.
\]
Subtracting stabilizers gives the Du Val type
\[
  2A_5+2A_2+2A_1,
\]
and the local motivic factor is
\[
  \bigl(V^{\mathrm{mot}}_{A_5}\bigr)^2
  \bigl(V^{\mathrm{mot}}_{A_2}\bigr)^2
  \bigl(V^{\mathrm{mot}}_{A_1}\bigr)^2 .
\]

\paragraph{Reduced motivic class.}
Maulik-Toda's reduced obstruction theory for a K3-fibred CY$_3$ uses
the K3 holomorphic two-form to define a cosection of the obstruction
sheaf. Since $g_N$ is symplectic, this two-form is $G_N$-invariant, and
the cosection descends through the finite \'etale quotient. With a
$G_N$-equivariant moduli problem and compatible orientation data fixed,
the reduced motivic class lies in $\KSring[\Lmot]$.

At $N=1$, the Maulik-Toda and Oberdieck-Shen reduced K3$\times E$
theorem gives the reduced DT/PT scalar. For $N\in\{2,3,4,6\}$, the CHL
motivic formula requires three additional hypotheses:
\[
\begin{gathered}
  \text{equivariant reduced motivic degeneration for the K3 fibration,}\\
  \text{descent of the orientation line through }G_N,\\
  \text{identification of the Kontsevich-Soibelman motivic}\\
  \text{quantum dilogarithm with the CHL denominator.}
\end{gathered}
\]
Under these hypotheses the identity is
\[
  \DTmot^{X_N}(q,p,y;\Lmot)
  =
  -\,\frac{C^{\mathrm{mot}}_N(q,y;\Lmot)}
          {(\FCHL_N(q,p,y))^{2,\mathrm{mot}}},
  \qquad
  \mathrm{wt}(\FCHL_N)=\frac{c_N(0)}2 .
\]

\paragraph{Primitive CHL Borcherds constants.}
The primitive CHL ladder is the five-level family
$N\in\{1,2,3,4,6\}$. Its weak Jacobi input is normalized by
\[
  c_N(0):=[q^0\zeta^0]\,\phi^{(g_N)}_{0,1}(\tau,z),\qquad
  \phi^{(\mathrm{id})}_{0,1}(\tau,z)=\zeta^{-1}+10+\zeta+O(q).
\]
Borcherds' weight formula gives the automorphic weight as $c_N(0)/2$.
For the Gritsenko-Nikulin and Gritsenko CHL denominator forms,
\[
\begin{array}{c|ccccc}
N&1&2&3&4&6\\\hline
c_N(0)&10&8&6&4&2\\
\kBKM(\FCHL_N)=c_N(0)/2&5&4&3&2&1
\end{array}
\]
The named denominator forms may be written
\[
  \Delta_5,\quad \Delta^{(2)}_4,\quad
  \Delta^{(3)}_3,\quad \Delta^{(4)}_2,\quad \Delta^{(6)}_1 .
\]
Thus the level-one CHL row has \(\kBKM(\Delta_5)=5\).

\paragraph{Gritsenko-Clery diagonal-divisor catalogue.}
The Gritsenko-Clery dd-modular catalogue is an eight-row atlas of
diagonal-divisor paramodular products. Its row labels and weights are
\[
\begin{array}{c|cccccccc}
(N,M)&(1,1)&(2,1)&(1,2)&(3,1)&(1,3)&(4,1)&(1,4)&(2,2)\\\hline
\mathrm{wt}(\FDD_{N,M})
&5&2&3&1&2&\frac12&\frac32&1\\
2\,\mathrm{wt}(\FDD_{N,M})
&10&4&6&2&4&1&3&2
\end{array}
\]
This catalogue has its own level, divisor, and multiplier data. It
shares the \((1,1)\) row with the primitive CHL ladder through
\(\Delta_5\). The other rows require their own comparison theorem before
their weights can be used in a CHL denominator statement.

\paragraph{The OP normalization.}
At $N=1$, the theta-normalized Gritsenko-Nikulin form has leading
coefficient $64$:
\[
  [q^{1/2}r^{1/2}s^{1/2}]\Delta_5=64,\qquad
  D_5=64^{-1}\Delta_5 .
\]
The unnormalized square and the Oberdieck-Pandharipande monic scalar
square are
\[
  \Phiun=\Delta_5^2,\qquad
  \PhiOP=D_5^2=4096^{-1}\Phiun .
\]
Consequently
\[
  Z_{\mathrm{OP/DT}}=-(\PhiOP)^{-1}=-D_5^{-2}
  =-4096\,\Delta_5^{-2}.
\]
The primitive BKM denominator is the square root \(D_5(2Z)\) in the
Gritsenko-Nikulin denominator convention. Squaring to \(\PhiOP\)
changes the scalar DT partition function and doubles the scalar weight.
The primitive BKM algebra remains the level-one \(\Delta_5\) algebra
with \(\kBKM=5\).

\paragraph{Wall crossing.}
Kontsevich-Soibelman motivic wall crossing applies to the equivariant
derived category after the orientation line and stability data are
chosen. The sector decomposition has the form
\[
  \MOmot(X_N;\gamma)
  =
  \MOmot(S\times E;\gamma)^{g_N\mathrm{-inv}}
  +\sum_{s=1}^{N-1}
   \MOmot(S\times E;\gamma)^{g_N,\zeta^s}\Lmot^{\mathrm{age}(s)} .
\]
This is an identity of motivic BPS classes in the motivic quantum torus.
Its numerical specialization is a character identity. A compact
Hall/BKM comparison is a stronger datum: it requires extension
correspondences, protected primitive bases, orientation signs, a
non-degenerate Hall pairing, a radical quotient, PBW compatibility, and
a comparison of brackets.

\paragraph{K-theoretic comparison.}
Oberdieck-Shen prove the reduced K-theoretic PT formula for
$S\times E$. The CHL extension requires the same reduced
$E$-fibred multiplicative structure after splitting every K-theoretic
class into $g_N$-invariant and twisted sectors. The twisted-sector
character is the Mathieu twining Jacobi form. The unconditional chain is
\[
  S\times E\ \mathrm{reduced\ DT}
  \longrightarrow
  S\times E\ K\text{-theoretic }PT
  \longrightarrow
  S\times E\ \mathrm{numerical\ DT}.
\]
For $N\in\{2,3,4,6\}$ the CHL chain is the same statement plus the
twined Oberdieck-Shen compatibility above.

\paragraph{Four invariants.}
The finite \'etale quotient gives
\[
  H^q(X_N,\mathcal O_{X_N})
  =
  H^q(S\times E,\mathcal O_{S\times E})^{G_N}.
\]
The symplectic K3 automorphism acts trivially on
$H^0(S,\mathcal O_S)$ and $H^2(S,\mathcal O_S)$, while the elliptic
translation acts trivially on $H^0(E,\mathcal O_E)$ and
$H^1(E,\mathcal O_E)$. Hence $h^{0,q}(X_N)=(1,1,1,1)$ and
\[
  \kcat(X_N)=\chi(\mathcal O_{X_N})=1-1+1-1=0 .
\]
In particular,
\[
  \kcat(K3\times E)=\chi(\mathcal O_{K3})\chi(\mathcal O_E)=2\cdot0=0.
\]
The compact CY$_3$ Hodge supertrace gives
\[
  \kch(X_N)=\sum_{q=0}^{3}(-1)^q h^{0,q}(X_N)=0 .
\]
For \(K3\times E\) the compact and chiral readings used in Vol III are
therefore the separated tuple
\[
\begin{array}{c|c|c}
\mathrm{construction}&\mathrm{invariant}&\mathrm{value}\\\hline
\mathrm{compact\ total\ space}&\kcat(K3\times E)&0\\
\mathrm{compact\ Hodge\ supertrace}&\kch(K3\times E)&0\\
\mathrm{Heisenberg\ specialization}&\kch^{\mathrm{Heis}}(K3\times E)&3\\
\mathrm{BKM\ denominator}&\kBKM(\Delta_5)&5\\
\mathrm{K3\ fibre\ Mukai\ rank}&\kfib(K3)&24
\end{array}
\]
The CHL Borcherds invariant is read from the Jacobi constant term:
\[
  \kBKM(\FCHL_N)=c_N(0)/2 .
\]
It is independent of \(\kcat(X_N)\), of the compact
\(\kch(X_N)\), and of the K3 fibre witness \(\kfib(K3)=24\). At
$N=1$, the Borcherds side has \(\kBKM(\Delta_5)=5\), while the compact
Hodge supertrace has \(\kch(K3\times E)=0\). At $N=2$, the Borcherds
side has \(\kBKM(\FCHL_2)=4\), while \(\kcat(X_2)=0\).

\paragraph{Statement.}
For the CHL quotient \(X_N\), the finite \'etale geometry gives a smooth
CY$_3$ source, the K3 holomorphic two-form gives the reduced cosection,
the primitive denominator gives
\(\kBKM(\FCHL_N)=c_N(0)/2\), and the compact total space gives
\(\kcat(X_N)=\kch(X_N)=0\). For \(N\in\{2,3,4,6\}\), equivariant
motivic degeneration, orientation descent, CHL motivic denominator
identification, and the twined Oberdieck-Shen K-theoretic comparison
give the motivic CHL identity displayed above.
